% Exact isolated candidate proof for source-aware verification.
% This file makes no novelty or priority claim and is not a live-book input.

\begin{proof}
\medskip\noindent\emph{Imported source scope.}
In \cite[Example~6.7.2]{CLD25}, Chambert--Loir and Ducros equip an
interval \(I\subseteq\overline{\mathbf R}\) with a G-tropical affine
structure for which parlinear functions are constant near every endpoint
belonging to the interval.  On a connected domain meeting such an
endpoint, every affine function is constant.  Sections~6.1.1 and 6.1.6
of \cite{CLD25} make tropical charts affine and their scalar coordinate
functions affine.  Together with \cite[Definition~6.2.1]{CLD25}, this
gives exactly \eqref{eq:p071-infinite-collar-smoothness} near the endpoint
\(-\infty\); away from that endpoint one uses the ordinary real affine
coordinate.  Conversely, an ordinary smooth function on the finite
stratum that is exactly constant on a terminal tail is represented by
the constant chart on that tail and by the ordinary coordinate on
endpoint-free neighbourhoods.  On the finite stratum,
\[
 \mathrm d'\mathrm d''f
       =f''(t)\,\mathrm d't\wedge\mathrm d''t,
\]
and the form vanishes on the terminal tail.  Hence positivity is exactly
\eqref{eq:p071-infinite-stage-zero}.  Definition~8.1.1 and
Lemma~8.1.4 of \cite{CLD25} give the generated sheaf and ordinal
filtration used in the statement.

\medskip\noindent\emph{Convexity and the endpoint trace.}
We first record two elementary consequences of the definition of
\(\mathscr J_{-\infty}\).  A locally extended-real convex function on a
connected ordinary interval is either finite everywhere or identically
\(-\infty\).  Indeed, the set where its value is \(-\infty\) is open:
if \(u(a)=-\infty\), express each nearby point as a strict convex
combination of \(a\) and a point on its other side.  The finite locus is
also open, since a finite value lying strictly between a nearby
\(-\infty\) value and a point on the other side would violate convexity.
The two loci are disjoint and cover a connected interval.

Suppose now that \(u\) is finite and convex on a tail
\((-\infty,b)\) and that
\begin{equation}
             c=\lim_{t\to-\infty}u(t)
                    \in\mathbf R\cup\{-\infty\}.
 \label{eq:p071-infinite-finite-side-limit}
\end{equation}
Then \(u\) is nondecreasing.  Otherwise there would be \(a<d<b\) with
\[
                 q=\frac{u(d)-u(a)}{d-a}<0.
\]
For every \(x<a\), monotonicity of secant slopes for a convex function
would give
\[
       \frac{u(a)-u(x)}{a-x}\leq q,
       \qquad
       u(x)\geq u(a)-q(a-x)\longrightarrow+\infty
                    \quad(x\to-\infty),
\]
contrary to \eqref{eq:p071-infinite-finite-side-limit}.  The same
nondecreasing conclusion is automatic when \(u\equiv-\infty\).

These observations also prove that \(\mathscr J_{-\infty}\) is a sheaf.
The ordinary convexity inequalities glue by the one-dimensional
local-to-global criterion.  Explicitly, on a compact interval take a
Lebesgue number for a finite cover by intervals of local convexity, choose
a partition whose consecutive triples lie in members of the cover, and
use the monotonicity of the resulting successive secant slopes.  The
dichotomy just proved handles extended-real values.  Finally, a member of
an open cover that contains \(-\infty\) contains a full terminal tail, so
its local equality \eqref{eq:p071-infinite-endpoint-trace} is already the
endpoint equality for the glued section.

The sheaf \(\mathscr J_{-\infty}\) is stable under pointwise decreasing
limits.  To see the only nonformal point, let \(u_j\downarrow u\) on an
endpoint component and put \(c_j=u_j(-\infty)\) and
\(c=u(-\infty)=\inf_jc_j\).  Every \(u_j\), and hence \(u\), is
nondecreasing, so
\begin{equation}
 c\leq L:=\lim_{t\to-\infty}u(t).
 \label{eq:p071-infinite-one-trace-inequality}
\end{equation}
If \(c\in\mathbf R\), then for every \(\varepsilon>0\) one can choose
\(j\) with \(c_j<c+\varepsilon\).  By the endpoint trace of \(u_j\), for
all sufficiently negative \(t\) one has
\(u(t)\leq u_j(t)<c+2\varepsilon\), so \(L\leq c\).  If
\(c=-\infty\), then, given \(M\in\mathbf R\), choose \(j\) with
\(c_j<M-1\) and then a sufficiently negative \(t\) so that
\(u(t)\leq u_j(t)<M\).  Monotonicity makes the same estimate valid farther
to the left, whence \(L=-\infty\).  Together with
\eqref{eq:p071-infinite-one-trace-inequality}, this proves preservation of
\eqref{eq:p071-infinite-endpoint-trace}.  Convexity passes directly to a
decreasing limit because no \(+\infty\) values occur.

Every member of \(\operatorname{Psh}_{S,0}\) lies in
\(\mathscr J_{-\infty}\): it is convex on the finite stratum, is constant
on a terminal tail, and has the required endpoint trace.  Since the latter
is a sheaf stable under decreasing limits, minimality gives
\begin{equation}
                 \operatorname{Psh}_S
                    \subseteq\mathscr J_{-\infty}.
 \label{eq:p071-infinite-generated-in-trace-class}
\end{equation}

\medskip\noindent\emph{Heat regularization with an exact terminal collar.}
Let \(-\infty\in W\), let \(u\in\mathscr J_{-\infty}(W)\), and shrink
the endpoint component so that
\([-\infty,R)\subseteq W\) for some \(R\in\mathbf R\).  If
\(u\equiv-\infty\) on the finite stratum of this tail, then the endpoint
trace gives \(u(-\infty)=-\infty\), and the smooth psh constants \(-n\)
give \eqref{eq:p071-infinite-fixed-domain-approximation} on every smaller
fixed tail.

We may therefore suppose that \(u\) is finite and convex on
\((-\infty,R)\).  Put \(c=u(-\infty)\).  By
\eqref{eq:p071-infinite-endpoint-trace} and the preceding argument, \(u\)
is nondecreasing.  Choose
\[
                       r<R_1<R_2<R
\]
and define
\[
              m=\frac{u(R_2)-u(R_1)}{R_2-R_1}\geq0.
\]
If \(c\in\mathbf R\), finite convexity gives continuity on
\((-\infty,R)\), while the trace equality gives continuity at the
endpoint \(-\infty\) in the order topology.  This proves the continuity
assertion in the statement.

Extend \(u\) to a finite function \(\bar u:\mathbf R\to\mathbf R\) by
\begin{equation}
 \bar u(t)=
 \begin{cases}
     u(t),&t\leq R_1,\\
     u(R_1)+m(t-R_1),&t\geq R_1.
 \end{cases}
 \label{eq:p071-infinite-convex-extension}
\end{equation}
The monotonicity of secant slopes shows that this extension is convex and
nondecreasing.  Every one of its secant slopes belongs to \([0,m]\), so
\begin{equation}
              |\bar u(x)-\bar u(y)|\leq m|x-y|
                    \quad(x,y\in\mathbf R).
 \label{eq:p071-infinite-global-lipschitz}
\end{equation}

For \(s>0\), set
\begin{equation}
 \gamma_s(y)=\frac{1}{\sqrt{4\pi s}}
       \mathrm e^{-y^2/(4s)},
 \qquad
 A_s(t)=\int_{\mathbf R}\bar u(t-y)\gamma_s(y)\,\mathrm dy.
 \label{eq:p071-infinite-heat-regularization}
\end{equation}
The linear-growth bound in
\eqref{eq:p071-infinite-global-lipschitz} makes the integral finite.
Every \(A_s\) is smooth, convex, and nondecreasing.  The Gaussian
semigroup identity and Jensen's inequality show that
\begin{equation}
             0<s'<s\quad\Longrightarrow\quad A_{s'}\leq A_s.
 \label{eq:p071-infinite-heat-monotonicity}
\end{equation}
Indeed, \(A_s=P_{s-s'}A_{s'}\), and averaging a convex function against
a centred probability measure can only increase it.  Moreover,
\begin{equation}
 0\leq A_s(t)-\bar u(t)
 \leq m\int_{\mathbf R}|y|\gamma_s(y)\,\mathrm dy
 =\frac{2m}{\sqrt{\pi}}\sqrt{s}.
 \label{eq:p071-infinite-heat-error}
\end{equation}

Choose a nonnegative even
\(\theta\in C^\infty_c((-1,1))\) of integral one.  For \(\eta>0\), put
\(\theta_\eta(s)=\eta^{-1}\theta(s/\eta)\) and
\begin{equation}
 M_\eta(a,b)=\int_{\mathbf R^2}
    \max\{a+s_1,b+s_2\}\theta_\eta(s_1)\theta_\eta(s_2)
       \,\mathrm ds_1\,\mathrm ds_2.
 \label{eq:p071-infinite-exact-maximum}
\end{equation}
This is a smooth convex function, nondecreasing in each variable, and
\begin{align}
 M_\eta(a,b)&=a &&\text{if }a>b+2\eta,
 \notag\\
 M_\eta(a,b)&=b &&\text{if }b>a+2\eta,
 \label{eq:p071-infinite-exact-maximum-regions}\\
 \max\{a,b\}&\leq M_\eta(a,b)\leq\max\{a,b\}+\eta.
 \notag
\end{align}
Fix \(\eta>0\) small enough that \(r+3\eta<R_1\), and define, for
\(n\geq1\),
\begin{equation}
 H_n(t)=M_\eta(t,-n)\quad(t\in\mathbf R),
 \qquad H_n(-\infty)=-n.
 \label{eq:p071-infinite-terminal-clamp}
\end{equation}
Then \(H_n\) is smooth, convex, and nondecreasing on the finite stratum,
is exactly equal to \(-n\) for \(t<-n-2\eta\), and is exactly equal to
\(t\) for \(t>-n+2\eta\).  Since \(M_\eta\) is nondecreasing in its
second variable,
\begin{equation}
                         H_{n+1}\leq H_n.
 \label{eq:p071-infinite-clamp-monotonicity}
\end{equation}

On the fixed open tail \(V=[-\infty,r)\), set
\begin{equation}
 F_n(t)=A_{1/n}(H_n(t))\quad(t\in V\cap\mathbf R),
 \qquad
 F_n(-\infty)=A_{1/n}(-n).
 \label{eq:p071-infinite-approximants}
\end{equation}
The function \(F_n\) is exactly constant on the endpoint collar
\(V\cap[-\infty,-n-2\eta)\).  On the finite stratum it is smooth and
\[
 F_n''
  =A_{1/n}''(H_n)(H_n')^2+A_{1/n}'(H_n)H_n''\geq0.
\]
Thus \(F_n\in\operatorname{Psh}_{S,0}(V)\).  By
\eqref{eq:p071-infinite-heat-monotonicity},
\eqref{eq:p071-infinite-clamp-monotonicity}, and monotonicity of
\(A_{1/n}\),
\begin{align}
 F_{n+1}
 &=A_{1/(n+1)}\circ H_{n+1}
 \leq A_{1/n}\circ H_{n+1}
 \leq A_{1/n}\circ H_n=F_n.
 \label{eq:p071-infinite-approximants-decrease}
\end{align}

For every finite \(t<r\), one has \(H_n(t)=t\) for all sufficiently large
\(n\).  Equation \eqref{eq:p071-infinite-heat-error} therefore gives
\(F_n(t)\to u(t)\).  For all sufficiently large \(n\), the point \(-n\)
lies in \((-\infty,R_1]\), so \(\bar u(-n)=u(-n)\).  At the endpoint, the
same estimate then gives
\begin{equation}
 |F_n(-\infty)-u(-n)|
   =|A_{1/n}(-n)-\bar u(-n)|
   \leq\frac{2m}{\sqrt{\pi n}},
 \label{eq:p071-infinite-endpoint-convergence}
\end{equation}
while \(u(-n)\to c=u(-\infty)\) in
\(\mathbf R\cup\{-\infty\}\).  This proves pointwise convergence on
the fixed tail, including when \(c=-\infty\).

If \(c\in\mathbf R\), the convergence is uniform on \(V\).  Indeed, for
all large \(n\), \(\bar u=u\) on the range \(H_n(V)\).  When
\(t>-n+2\eta\), the clamp is the identity.  When
\(t\leq-n+2\eta\), the bounds in
\eqref{eq:p071-infinite-exact-maximum-regions} give
\[
                 t\leq H_n(t)\leq-n+3\eta,
\]
and monotonicity gives
\[
 0\leq u(H_n(t))-u(t)\leq u(-n+3\eta)-c.
\]
For finite \(t\), the inequalities \(H_n(t)\geq t\),
\(A_{1/n}\geq\bar u\), and monotonicity of \(\bar u\) also give
\(F_n(t)\geq u(t)\).  At the endpoint,
\[
 0\leq F_n(-\infty)-c
 \leq\frac{2m}{\sqrt{\pi n}}+u(-n)-c
 \leq\frac{2m}{\sqrt{\pi n}}+u(-n+3\eta)-c.
\]
Together with \eqref{eq:p071-infinite-heat-error}, these estimates yield
\[
 0\leq\sup_{t\in V}\bigl(F_n(t)-u(t)\bigr)
 \leq\frac{2m}{\sqrt{\pi n}}
       +u(-n+3\eta)-c\longrightarrow0.
\]

\medskip\noindent\emph{Interior points and conclusion.}
At an interior point, the extended-real dichotomy reduces matters either
to the constants \(-n\), or to a finite convex function on an ordinary
interval.  In the latter case, restrict to a compact subinterval and
extend the function to a finite convex function on \(\mathbf R\) by
affine supporting lines at its two endpoints.  Its heat regularizations
\(A_{1/n}\) are smooth and convex, decrease by
\eqref{eq:p071-infinite-heat-monotonicity}, and converge to the original
function.  Thus every section of \(\mathscr J_{-\infty}\) belongs locally
to \(\operatorname{Psh}_{S,1}\).  Hence
\[
 \mathscr J_{-\infty}\subseteq\operatorname{Psh}_{S,1}
 \subseteq\operatorname{Psh}_S.
\]
Combining this with
\eqref{eq:p071-infinite-generated-in-trace-class} proves
\eqref{eq:p071-infinite-first-stage-collapse} and all the stated
approximation assertions.  Finally, \(u(t)=t\), with endpoint value
\(-\infty\), satisfies the trace condition and is obtained explicitly
by the same construction.
\end{proof}
